\section{Priprava za prenos v realno okolje}
Cilj tega prispevka ni simulacija sama po sebi ampak mož\-nost prenosljivosti predlagane rešitve na realni sistem. Največji problem delovanja Sim2Real (angl. sim-to-real) je nepopolnost realnega okolja in strojne opreme, kar ni vedno mogoče doseči z virtualnim modelom. To rešimo z integracijo rekurzivne arhitekture GRU in strogim upo\-šte\-va\-njem omejitev strojne opreme. Upoštevali smo velikost pomnilnika mikrokrmilnika STM32F446RE, zaradi česar smo nevronsko mrežo omejili na le tri skrite sloje z 32 nevroni oz. GRU celicami (slika \ref{nevronska mreža}). Nevronska mreža s skritim slojem 32 GRU celic in dvema skritima slojema po 32 nevronov vsebuje približno 6600 parametrov, kar v 32-bitnem zapisu s plavajočo vejico znaša okoli 26 kB. To predstavlja zanemarljiv delež pomnilnika mikrokrmilnika STM32F446RE (512 kB bliskovnega pomnilnika in 128 kB SRAM), rezultati sklepanja pa zahtevajo manj kot 1 kB delovnega pomnilnika. En prehod skozi mrežo obsega več tisoč množenj in seštevanj, kar na jedru Cortex-M4 traja nekaj $10\,\mu\text{s}$, kar je bistveno manj od regulacijskega koraka 20 ms. Za čim večjo računsko hitrost smo uporabili fiksno normalizacijo vhodov nevronske mreže s konstantami, za aktivacijsko funkcijo pa izbrali računsko nezahtevno funkcijo ReLU. Na realnem sistemu kontrolna zanka zaradi omejitev komunikacijskih vodil CAN in RS\-485 teče s frekvenco 50 Hz, zato smo regulacijo z nevronsko mrežo na enako frekvenco omejili tudi v simulaciji. Zaradi zaščite motorjev DDSM115 smo prav tako omejili izhod nevronske mreže za tok kolesa s funkcijo $\tanh$, ki tok zvezno preslika v območje $\pm 2{,}0$~A in tako prepreči preseganje varne tokovne meje.
  
  Sim2Real smo želeli izboljšati še z dodajanjem randomizacije domene (angl. domain randomization), pri čemer celice GRU skozi sprotno posodabljanje skritega stanja delujejo kot implicitni identifikator spreminjajočega se sistema. Namesto da bi poskušali parametre realne platforme natančno izmeriti in poustvariti v simulaciji, strategijo učimo na celotni porazdelitvi verjetnih vrednosti, s čimer predvidevamo, da je velika možnost, da realna platforma deluje z vrednostmi, ki padejo znotraj te porazdelitve. Izdelali smo funkcijo, ki na začetku vsake epizode vpelje randomizacijo domene in naključno spreminja ojačenje motorjev (od $0{,}8$ do $1{,}2$), prag prenehanja delovanja motorjev zaradi trenja (od $0{,}03$ do $0{,}20$~A), odmik toka zaradi nesimetrije motorja (od $-0{,}08$ do $0{,}08$~A), časovno konstanto tokovne zanke (od $0{,}005$ do $0{,}020$~s), tokovno omejitev posameznega motorja (od $1{,}6$ do $2{,}4$~A), odmik meritve naklona zaradi napake v montaži ali kalibraciji IMU senzorja (od $-0{,}5$ do $0{,}5$°) ter zakasnitev akcij na motorjih in zakasnitev opazovanj senzorjev (naključno $0$ ali $1$ korak). Prisotnost zakasnitev in nelinearnih šumov (Gaussov šum naklona s $\sigma = 0{,}15$° in kotne hitrosti s $\sigma = 0{,}02$~rad/s) v vsakem koraku simulacije dodatno utemeljuje uporabo arhitekture GRU, saj mreža s spominom uspešno kompenzira delno opazovanje stanja brez vnosa faznega zamika.